\documentclass[10pt,a4paper]{amsart}
\usepackage[margin=19mm]{geometry}
\usepackage{amsmath,amssymb,mathtools}
\usepackage[scr=rsfs,cal=euler]{mathalfa}
\usepackage{xfrac}
\usepackage[unicode,hidelinks,bookmarks=false]{hyperref}
\newcommand{\ie}{i.e.\ }
\newcommand{\cf}{cf.\ }
\newcommand{\resp}{respectively\ }
\newcommand{\Subsection}{\subsection}
\providecommand{\nobreakdash}{\nobreak\hbox{-}}
\setlength{\emergencystretch}{2em}
\multlinegap0pt
\usepackage{bm}
\usepackage{bbm}
\usepackage{dsfont}
\usepackage{xspace}
\usepackage{cancel}
\usepackage{mathtools}
\usepackage{amsthm}
\makeatletter
\@ifundefined{lemma}{\newtheorem{lemma}{Lemma}}{}
\makeatother
\title[Test-Time Steering for Lossless Text Compression via Weighte]{Test-Time Steering for Lossless Text Compression via Weighted Product of Experts}
\author{Qihang Zhang, Muchen Li, Ziao Wang, Renjie Liao,, Lele Wang}
\date{Source: arXiv:2511.10660v1 (CC BY 4.0)}
\begin{document}
\maketitle

\noindent\emph{Source and licence.} Test-Time Steering for Lossless Text Compression via Weighted Product of Experts. Version arXiv:2511.10660v1 (CC BY 4.0), distributed under the Creative Commons Attribution 4.0 International licence, as recorded in the arXiv licence metadata for this exact version and shown on the abstract page https://arxiv.org/abs/2511.10660v1. Full rights notice: licenses/arxiv-2511.10660-CC-BY-4.0.txt.\par\smallskip

\noindent\textit{Editorial scope.} Counted unit: the Lemma whose source label is \texttt{lemma:proof\_inequality} in the article (source0.tex lines 595--608, source label \texttt{lemma:proof\_inequality}), delivered together with the article's own complete original proof of it (source0.tex lines 685--763, one \texttt{proof} environment of 79 lines). No mathematics is written, repaired or re-derived: statement and proof are the article's own text line for line. The proof is detached from the statement in the source: the article states the result at lines 595--608 and prints its proof at lines 685--763, and the proof environment's own header names the statement, so the association is the article's own. Required context retained uncounted: none. Every definition, display and label of those lines is retained unchanged. Omitted from this package and not redistributed: 1--594, 609--684, 764--844. Nothing inside a retained excerpt is omitted or altered: every retained body is delivered exactly as the article's lines are written, so no omission receipt of any kind accompanies this package. Citations: the retained excerpts carry no citation command at all, so no bibliography entry of the article is transcribed and no reference is redistributed. Reference adaptation: none was needed and none was made; every reference inside the delivered unit resolves inside it, so no unresolved reference or citation is produced. Printed slips kept literal: no printed word, symbol, hypothesis or number of the counted statement or of its proof was altered. Layout and class: the article's own class is \texttt{article} with packages including acl, times, latexsym, fontenc, inputenc, microtype, inconsolata, graphicx, booktabs, tabularx, amsmath, amsfonts, mathtools, hyperref; the wrapper is the lane's pinned \texttt{amsart} editorial wrapper at 10pt on A4, which restates the theorem-like environments the retained text needs and adds the article's own macro definitions that the retained text needs (none), with extra wrapper packages (bm, bbm, dsfont, xspace, cancel, mathtools). The delivered numbering is the unit's own. Assets: no retained span contains a figure environment, a graphics-inclusion command, a PDF-page inclusion, a table or a TikZ picture; no image file is copied into this package, the package declares no asset dependency and no image file is redistributed with it. Licence: version arXiv:2511.10660v1 of this article is distributed under the Creative Commons Attribution 4.0 International licence, as recorded in the arXiv licence metadata for this exact version and shown on the abstract page https://arxiv.org/abs/2511.10660v1. A full rights notice is delivered at licenses/arxiv-2511.10660-CC-BY-4.0.txt. Characters: every retained excerpt is pure ASCII, so the delivered engine, XeLaTeX with its default Latin Modern text font, reports no missing character.
\begin{lemma}
    \label{lemma:proof_inequality}
    Let $p^{(k)} = \bigl(p^{(k)}_1, \ldots, p^{(k)}_D\bigr)$ for $k=1,\dots,K$ 
    be $K$ categorical distributions, so $\sum_{j=1}^D p^{(k)}_j = 1$ for each $k$. 
    Let $\alpha_1,\dots,\alpha_K \ge 0$ satisfy $\sum_{k=1}^K \alpha_k = 1.$ 
    Then
    \[
    \sum_{j=1}^D 
    \prod_{k=1}^K \bigl(p^{(k)}_j\bigr)^{\alpha_k}
    \;\;\le\;\; 1,
    \]
    with equality if and only if $p^{(1)} = p^{(2)} = \cdots = p^{(K)}$ 
    or exactly one $\alpha_k=1$ and the rest are zero.
\end{lemma}

\begin{proof}[Proof of Lemma~\ref{lemma:proof_inequality}]

\textbf{Base Case ($K=2$).} 
For two distributions $(p_j)$ and $(q_j)$ with weights $\alpha$ and $1-\alpha$, 
it is already proved that
\[
\sum_{j=1}^D p_j^{\,\alpha} \, q_j^{\,1-\alpha} \;\;\le\;\; 1,
\]
with equality precisely if $p=q$ or $\alpha \in \{0,1\}$.

\textbf{Inductive Step.} 
Assume the statement holds for $K-1$ distributions. We prove it for $K$.

Let $p^{(1)}, p^{(2)}, \dots, p^{(K)}$ be $K$ distributions with weights 
$\alpha_1,\dots,\alpha_K$ such that $\sum_{k=1}^K \alpha_k = 1$. 
Define $\alpha' = \alpha_1 + \alpha_2$. 
By the $K=2$ base case,
\[
\sum_{j=1}^D 
\bigl(p^{(1)}_j\bigr)^{\frac{\alpha_1}{\alpha'}}
\bigl(p^{(2)}_j\bigr)^{\frac{\alpha_2}{\alpha'}}
\;\;\le\;\; 1.
\]
Given this sum can not be zero, define a new ``composite'' distribution 
$r = (r_1,\dots,r_D)$ by
\[
r_j 
= 
\frac{
 \bigl(p^{(1)}_j\bigr)^{\frac{\alpha_1}{\alpha'}}\,\bigl(p^{(2)}_j\bigr)^{\frac{\alpha_2}{\alpha'}}
}{
 \sum_{i=1}^D 
 \bigl(p^{(1)}_i\bigr)^{\frac{\alpha_1}{\alpha'}}\,\bigl(p^{(2)}_i\bigr)^{\frac{\alpha_2}{\alpha'}}
},
\quad
j=1,\dots,D.
\]
Clearly, $\sum_{j=1}^D r_j = 1$. 
Now we have $K-1$ distributions: $r, p^{(3)}, \dots, p^{(K)}$, 
with weights $\alpha',\alpha_3,\dots,\alpha_K$ summing to 1.

By the inductive hypothesis, 
\[
\sum_{j=1}^D
r_j^{\,\alpha'}
\,\bigl(p^{(3)}_j\bigr)^{\alpha_3}
\,\cdots\,
\bigl(p^{(K)}_j\bigr)^{\alpha_K}
\;\;\le\;\; 1.
\]
Substituting the definition of $r_j$ into the left-hand side shows
\begin{multline}
\sum_{j=1}^D
  (p^{(1)}_j)^{\alpha_1}(p^{(2)}_j)^{\alpha_2}
  \prod_{k=3}^K (p^{(k)}_j)^{\alpha_k} \\
\le
\Biggl(
  \sum_{i=1}^D
    (p^{(1)}_i)^{\alpha_1}(p^{(2)}_i)^{\alpha_2}
\Biggr)^{\alpha'}.
\end{multline}
Since the base case says $\sum_{i=1}^D (p^{(1)}_i)^{\alpha_1} (p^{(2)}_i)^{\alpha_2} 
\le 1$, 
raising it to the power $\alpha'$ also yields a value at most 1. Thus
\[
\sum_{j=1}^D
\prod_{k=1}^K \bigl(p^{(k)}_j\bigr)^{\alpha_k}
\;\;\le\;\; 1,
\]
which completes the inductive step.

\textbf{Equality Condition.} 
In base case, equality demands $p=q$ or 
$\alpha=0 $ or $\alpha = 1$. Propagating this requirement through 
each step of the recursion shows all distributions must be identical 
or all but one weight are zero. 

By induction, the proof is complete.
\end{proof}

\end{document}
